% id: 103-16
% book: 베이직쎈 미적분1 · 06 도함수의 활용 (3)
% section: 유형 067 그래프에서의 속도와 가속도
% figure: tikz:fig-103-16
% answer: 
% answer_source: 
% variant_level: 0 (원본 전사)
% 이 파일은 build.mjs 가 items.json 에서 만든다. 고칠 때는 items.json 을 고친다.
\smash{\raisebox{1.5mm}{\dmkichul{베이직쎈 미적분1 103쪽 16번}}}
\noindent\begin{minipage}[t]{0.58\linewidth}\vspace{0pt}\raggedright 원점에서 출발하여 수직선 위를 움직이는 점~$\pt{P}$의 시각 $t$에서의 속도 $v(t)$의 그래프가 오른쪽 그림과 같을 때, 옳은 것만을 \textbf{보기}에서 있는 대로 고르시오. \end{minipage}\hfill\begin{minipage}[t]{0.4\linewidth}\vspace{0pt}\centering \resizebox{\ifdim\width>\linewidth\linewidth\else\width\fi}{!}{% 103-16(103쪽) — 원점에서 출발한 점 P 의 시각 t 에서의 속도 v(t) 의 그래프(0~1 증가, 1~3 일정, 3~5 감소해 t=4 에서 0, 5~7 증가해 t=6 에서 0, 7 뒤로 일정). 스캔 화질이 낮아 TikZ 로 다시 그림.
% 원문에 있는 것만 옮긴다: 꺾은선 · 안내 점선 네 줄(1, 3, 5, 7) · 라벨 v(t), t, O, 1, 3, 4, 5, 6, 7. 세로축에는 원문처럼 눈금값을 적지 않는다.
\begin{tikzpicture}[x=0.30cm, y=0.30cm, line width=0.6pt]
  \draw[->, >=stealth, line width=0.7pt] (-0.7,0) -- (9.3,0) node[below=1pt, font=\small] {$t$};
  \draw[->, >=stealth, line width=0.7pt] (0,-2.9) -- (0,3.0) node[left=1pt, font=\small] {$v(t)$};
  \draw[densely dashed, line width=0.4pt] (1,0) -- (1,2.2);
  \draw[densely dashed, line width=0.4pt] (3,0) -- (3,2.2);
  \draw[densely dashed, line width=0.4pt] (5,0) -- (5,-2.2);
  \draw[densely dashed, line width=0.4pt] (7,0) -- (7,2.2);
  \draw (0,0) -- (1,2.2) -- (3,2.2) -- (5,-2.2) -- (7,2.2) -- (8.5,2.2);
  \node[below left=-2pt and -2.5pt, font=\small] at (0,0) {$\pt{O}$};
  \node[below=1pt, font=\small] at (1,0) {$1$};
  \node[below=1pt, font=\small] at (3,0) {$3$};
  \node[below=1pt, font=\small] at (4,0) {$4$};
  \node[above=0.5pt, font=\small] at (5,0) {$5$};
  \node[below=1pt, font=\small] at (6,0) {$6$};
  \node[below=1pt, font=\small] at (7,0) {$7$};
\end{tikzpicture}
}\end{minipage}\par
\noindent \bogi{ㄱ. $1<t<3$에서 점~$\pt{P}$의 가속도는 일정하다.\\ ㄴ. $t=4$에서 점~$\pt{P}$의 가속도는 $0$이다.\\ ㄷ. $t=5$에서와 $t=7$에서의 점~$\pt{P}$의 운동 방향은 서로 반대이다.\\ ㄹ. $0<t<7$에서 점~$\pt{P}$는 운동 방향을 세 번 바꾼다.}
